\documentclass[./examen.tex]{subfiles}

\begin{document}
\begin{center}
\makebox[\textwidth]{Nombre y Apellido:\enspace\hrulefill}
\end{center}
\begin{multicols}{2}
\begin{questions}
 \question Un electroimán tiene una Fuerza de $F=10Kp$ y su sección es de $S=0.002m^2$.
 \begin{parts}
    \part[2] Calcular la induccción $B$ del electroimán en $Tesla$ y calcular la intensidad de campo magnético $H$ si la permeabilidad relativa del núcleo de electroimán es $\mu_r=1000$
 \end{parts}
\question El circuito de la figura está construído de chapa magnética $\mu_r=2000$, se desea una inducción de $B=2T$ la bobina es de $N=1500$ espiras.  
\begin{parts}
\part[2] Calcular la longitud de la línea media $l_m$ en metros. 
\part[2] Calcular la Intensidad de campo en el núcleo $H$ en $Av/m$.
\part[2] Fuerza magnetomotriz $\mathscr{F}$ amperios-vuelta necesarios.
\part[2] Intensidad de corriente que circula por la bobina.
\end{parts}
\end{questions}
\columnbreak
\centering
\begin{tikzpicture}[scale=0.035]
\draw [ultra thick] (0,0) rectangle (180,180);
\draw [ultra thick] (40,40) rectangle (140,140);
\draw [thin,dashed] (20,20) rectangle (160,160);
\draw [thin,Circle -] (-20,140)--(40,140);
\draw [thin] (0,130)--(40,130);
\draw [thin] (0,120)--(40,120);
\draw [thin] (0,110)--(40,110);
\draw [thin] (0,100)--(40,100);
\draw [thin] (0,90)--(40,90);
\draw [thin] (0,80)--(40,80);
\draw [thin] (0,70)--(40,70);
\draw [thin] (0,60)--(40,60);
\draw [thin] (0,50)--(40,50);
\draw [thin,Circle -] (-20,40)--(0,40);
\draw [stealth-stealth] (140,80)--(180,80);
\draw [stealth-stealth] (40,100)--(140,100);
\draw [stealth-stealth] (100,140)--(100,40);
\draw [stealth-stealth] (80,40)--(80,0);
\draw (80,110) node[]{5cm};
\draw (110,70) node[]{5cm};
\draw (170,90) node[]{2cm};
\draw (70,30) node[]{2cm};
\end{tikzpicture}

\end{multicols}




 \end{document}